\documentclass[11pt,a4paper]{article}

\usepackage[utf8]{inputenc}
\usepackage[T1]{fontenc}
\usepackage[french]{babel}
\usepackage[a4paper,margin=1.8cm]{geometry}
\usepackage{amsmath,amssymb,amsthm,mathtools}
\usepackage{lmodern}
\usepackage{xcolor}
\usepackage{array,tabularx,booktabs,multirow}
\usepackage{enumitem}
\usepackage{fancyhdr}
\usepackage{lastpage}
\usepackage{hyperref}
\usepackage{graphicx}
\usepackage{tikz}
\usepackage{pgfplots}
\usepackage{tcolorbox}
\usepackage{multicol}
\usepackage{siunitx}
\usetikzlibrary{arrows.meta,calc,positioning,patterns,decorations.pathmorphing}
\pgfplotsset{compat=1.18}
\hypersetup{colorlinks=true,linkcolor=blue,urlcolor=blue}
\sisetup{locale=FR}

\setlength{\headheight}{14pt}
\pagestyle{fancy}
\fancyhf{}
\fancyhead[L]{Première générale -- Spécialité Physique-Chimie}
\fancyhead[C]{Cours détaillé -- Ondes et signaux}
\fancyhead[R]{Page \thepage/\pageref{LastPage}}
\fancyfoot[C]{Pierre Garcia-Buscaylet -- support de cours LaTeX}

\definecolor{bleu1}{RGB}{30,90,160}
\definecolor{bleu2}{RGB}{225,238,255}
\definecolor{vert1}{RGB}{30,120,80}
\definecolor{vert2}{RGB}{231,247,239}
\definecolor{rouge1}{RGB}{165,55,55}
\definecolor{rouge2}{RGB}{252,235,235}
\definecolor{jaune2}{RGB}{255,248,220}
\definecolor{violet2}{RGB}{243,236,255}

\newtcolorbox{definitionbox}{colback=bleu2,colframe=bleu1,title=Définition,fonttitle=\bfseries}
\newtcolorbox{propriete}{colback=vert2,colframe=vert1,title=Propriété / résultat,fonttitle=\bfseries}
\newtcolorbox{methode}{colback=jaune2,colframe=orange!80!black,title=Méthode,fonttitle=\bfseries}
\newtcolorbox{attentionbox}{colback=rouge2,colframe=rouge1,title=Attention / piège,fonttitle=\bfseries}
\newtcolorbox{ouverture}{colback=violet2,colframe=purple!70!black,title=Ouverture Terminale,fonttitle=\bfseries}
\newtcolorbox{exemplebox}{colback=white,colframe=black!65,title=Exemple,fonttitle=\bfseries}

\newtheorem*{demonstration}{Démonstration}

\begin{document}

\begin{center}
    {\LARGE\bfseries Ondes et signaux}\\[0.3em]
    {\large Cours détaillé de Première spécialité Physique-Chimie}\\[0.4em]
    \rule{0.9\linewidth}{0.6pt}
\end{center}

\tableofcontents
\bigskip

\section*{Comment lire ce chapitre}
Ce document ne se contente pas d'énoncer des formules : il cherche à faire comprendre le \emph{sens physique} des notions. À chaque fois, on répond à quatre questions :
\begin{enumerate}[label=\arabic*.]
    \item \textbf{De quoi parle-t-on ?}
    \item \textbf{Que voit-on expérimentalement ?}
    \item \textbf{Quelle relation mathématique faut-il savoir utiliser ?}
    \item \textbf{Pourquoi cela prépare-t-il la Terminale ?}
\end{enumerate}

\section{Ondes mécaniques}

\subsection{Qu'est-ce qu'une onde mécanique progressive ?}

\begin{definitionbox}
Une \textbf{onde mécanique progressive} est la propagation d'une \textbf{perturbation} dans un \textbf{milieu matériel}, sans transport global de matière, mais avec transport d'énergie.
\end{definitionbox}

\paragraph{Idée essentielle.}
Quand on secoue une corde, ce n'est pas la corde entière qui se déplace jusqu'à l'autre bout : ce qui se propage, c'est la \emph{déformation}. De même, pour un son, ce sont des compressions et des dilatations de l'air qui se transmettent.

\begin{exemplebox}
Exemples d'ondes mécaniques :
\begin{itemize}
    \item vague à la surface de l'eau ;
    \item onde le long d'une corde ;
    \item onde sonore dans l'air ;
    \item onde sismique dans la Terre ;
    \item houle en mer.
\end{itemize}
Contre-exemple : la lumière n'est \textbf{pas} une onde mécanique, car elle peut se propager dans le vide.
\end{exemplebox}

\begin{center}
\begin{tikzpicture}[scale=1]
    \draw[thick,->] (0,0) -- (10,0) node[right] {direction de propagation};
    \draw[blue,very thick,domain=0:9,samples=120] plot(\x,{0.7*sin(deg(1.2*\x))});
    \draw[red,fill=red] (1.3,0.70) circle (0.06);
    \draw[red,->,thick] (1.3,0.70) -- ++(0,0.7) node[above] {mouvement local du milieu};
    \node at (5,-1.1) {Une onde sur une corde : la perturbation avance, mais la corde ne part pas avec l'onde.};
\end{tikzpicture}
\end{center}

\begin{propriete}
Une onde mécanique nécessite un milieu matériel : elle ne se propage ni dans le vide, ni sans support physique.
\end{propriete}

\subsection{Propagation d'une perturbation : interprétation qualitative}
Dans un milieu matériel, les particules sont couplées entre elles. Si l'une est déplacée, elle agit sur sa voisine, puis celle-ci sur la suivante, etc. C'est ce \textbf{couplage local} qui permet la propagation.

\begin{methode}
Pour expliquer qualitativement une onde mécanique, on doit dire :
\begin{enumerate}
    \item qu'une zone du milieu est perturbée ;
    \item que cette zone agit sur la zone voisine ;
    \item que l'effet se transmet de proche en proche ;
    \item qu'il y a propagation d'énergie mais pas transport de matière à grande échelle.
\end{enumerate}
\end{methode}

\begin{attentionbox}
Beaucoup d'élèves disent : \og l'onde déplace la matière \fg. Ce n'est en général pas le bon modèle. Les particules du milieu oscillent localement autour d'une position d'équilibre.
\end{attentionbox}

\subsection{Grandeurs associées à une onde}
Pour décrire une onde, on utilise notamment :
\begin{itemize}
    \item l'\textbf{amplitude} : importance maximale de la perturbation ;
    \item la \textbf{durée} de propagation ;
    \item la \textbf{distance} parcourue ;
    \item la \textbf{célérité} : vitesse de propagation de la perturbation dans le milieu.
\end{itemize}

\begin{definitionbox}
La \textbf{célérité} $c$ d'une onde est la vitesse à laquelle la perturbation se propage dans le milieu.
\[
 c = \frac{d}{\Delta t}
\]
où $d$ est la distance parcourue et $\Delta t$ la durée de propagation.
\end{definitionbox}

\subsection{Retard et durée de propagation}

\begin{definitionbox}
Le \textbf{retard} entre deux points $A$ et $B$ est le temps mis par la perturbation pour aller de $A$ à $B$.
\[
\tau = \frac{AB}{c}
\]
\end{definitionbox}

\begin{exemplebox}
Un son se propage dans l'air à environ \SI{340}{m.s^{-1}}. Si un éclair d'orage est vu, puis le tonnerre entendu \SI{4,0}{s} plus tard, on estime la distance de l'orage par
\[
 d = c\,\Delta t = 340\times 4,0 \approx \SI{1360}{m}.
\]
L'orage est donc à environ \SI{1,4}{km}.
\end{exemplebox}

\begin{demonstration}
La relation $c=\dfrac{d}{\Delta t}$ traduit exactement la définition d'une vitesse moyenne : distance parcourue divisée par la durée correspondante. Ici, cette vitesse n'est pas celle d'un objet matériel, mais celle de la perturbation.
\end{demonstration}

\subsection{Influence du milieu sur la célérité}
La célérité d'une onde dépend du milieu. Une même source peut produire des ondes de fréquences identiques dans plusieurs milieux, mais la célérité change.

\begin{propriete}
La célérité d'une onde mécanique dépend des propriétés du milieu : rigidité, inertie, température, densité, etc.
\end{propriete}

\begin{exemplebox}
\begin{itemize}
    \item Le son va plus vite dans l'eau que dans l'air.
    \item Une onde se propage plus vite sur une corde plus tendue.
    \item La lumière, elle, a une célérité particulière dans le vide, notée $c$, qui vaut environ $3,00\times 10^8\ \si{m.s^{-1}}$.
\end{itemize}
\end{exemplebox}

\begin{ouverture}
En Terminale, on rencontrera des phénomènes d'interférences, de diffraction et l'effet Doppler. Pour bien les comprendre, il faut déjà être très solide sur la notion de célérité, de fréquence et de longueur d'onde.
\end{ouverture}

\subsection{Onde périodique}

\begin{definitionbox}
Une onde est dite \textbf{périodique} lorsque la perturbation se répète identiquement à intervalles de temps réguliers.
\end{definitionbox}

\paragraph{Exemple mental.}
Si l'on agite régulièrement l'extrémité d'une corde de haut en bas, on crée un motif qui se reproduit. Ce motif se propage : on obtient une onde périodique.

\subsection{Périodicité temporelle et périodicité spatiale}

\begin{definitionbox}
\begin{itemize}
    \item La \textbf{période temporelle} $T$ est la durée au bout de laquelle, en un point donné, le signal se répète.
    \item La \textbf{longueur d'onde} $\lambda$ est la distance parcourue par l'onde pendant une période. C'est aussi la distance entre deux points du milieu qui vibrent en phase.
\end{itemize}
\end{definitionbox}

\begin{center}
\begin{tikzpicture}[scale=1]
    \draw[->] (0,0) -- (11,0) node[right] {$x$};
    \draw[->] (0,-1.7) -- (0,1.7) node[above] {élongation};
    \draw[blue,very thick,domain=0:10,samples=200] plot(\x,{sin(deg(2*pi*\x/4))});
    \draw[dashed] (1,0) -- (1,1);
    \draw[dashed] (5,0) -- (5,1);
    \draw[<->,red,thick] (1,-1.2) -- (5,-1.2);
    \node[red] at (3,-1.5) {$\lambda$};
    \node at (6.5,1.2) {profil spatial à un instant donné};
\end{tikzpicture}
\end{center}

\begin{propriete}
La relation fondamentale des ondes périodiques est
\[
\lambda = cT
\]
Comme $f=\dfrac{1}{T}$, on peut aussi écrire
\[
c = \lambda f.
\]
\end{propriete}

\begin{demonstration}
Pendant une période $T$, la perturbation avance d'une distance égale à une longueur d'onde $\lambda$. En appliquant la définition de la célérité,
\[
c=\frac{\text{distance parcourue}}{\text{durée}}=\frac{\lambda}{T},
\]
d'où
\[
\lambda=cT.
\]
Comme $f=1/T$, on obtient aussi $c=\lambda f$.
\end{demonstration}

\begin{attentionbox}
\textbf{Ne pas confondre :}
\begin{itemize}
    \item $T$ s'exprime en secondes : c'est une durée ;
    \item $\lambda$ s'exprime en mètres : c'est une distance.
\end{itemize}
L'une décrit la répétition dans le temps, l'autre dans l'espace.
\end{attentionbox}

\subsection{Onde sinusoïdale}
Une onde sinusoïdale est un cas particulier très important, car elle est simple mathématiquement et très fréquente en physique.

\begin{definitionbox}
Une onde sinusoïdale est une onde dont l'évolution temporelle ou spatiale peut être modélisée par une fonction sinus ou cosinus.
\end{definitionbox}

\paragraph{En un point fixe du milieu.}
On peut écrire, pour une grandeur $y(t)$,
\[
y(t)=A\sin\left(2\pi \frac{t}{T}+\varphi\right)
\]
où $A$ est l'amplitude et $\varphi$ une phase initiale.

\paragraph{Intérêt du modèle.}
Même si tous les signaux réels ne sont pas purement sinusoïdaux, le sinus sert de brique élémentaire pour comprendre les phénomènes périodiques.

\begin{ouverture}
En Terminale, le modèle sinusoïdal deviendra central pour les circuits électriques en régime alternatif, pour la diffraction, pour les interférences et pour les oscillateurs.
\end{ouverture}

\subsection{Lire un graphique temporel ou spatial}

\begin{methode}
\textbf{Sur un graphique temporel} :
\begin{itemize}
    \item on lit la période $T$ entre deux maxima successifs ;
    \item on lit l'amplitude comme l'écart maximal à la position d'équilibre.
\end{itemize}
\textbf{Sur un graphique spatial} :
\begin{itemize}
    \item on lit la longueur d'onde $\lambda$ entre deux crêtes successives ;
    \item si on connaît $T$ ou $f$, on en déduit $c$.
\end{itemize}
\end{methode}

\subsection{Exemples d'application}

\begin{exemplebox}
Une onde sonore a une fréquence $f=\SI{440}{Hz}$ dans l'air. On prend $c=\SI{340}{m.s^{-1}}$.
\[
\lambda=\frac{c}{f}=\frac{340}{440}\approx \SI{0,77}{m}.
\]
La longueur d'onde vaut environ \SI{77}{cm}.
\end{exemplebox}

\begin{exemplebox}
Une perturbation met \SI{0,25}{s} pour parcourir \SI{80}{cm} le long d'une corde.
\[
c=\frac{0,80}{0,25}=\SI{3,2}{m.s^{-1}}.
\]
\end{exemplebox}

\subsection{Synthèse sur les ondes mécaniques}
\begin{center}
\renewcommand{\arraystretch}{1.35}
\begin{tabularx}{\linewidth}{|>{\bfseries}m{4.1cm}|X|}
\hline
Onde mécanique progressive & Propagation d'une perturbation dans un milieu matériel. \\
\hline
Célérité $c$ & Vitesse de propagation de la perturbation. \\
\hline
Retard $\tau$ & Temps mis pour aller d'un point à un autre. \\
\hline
Période $T$ & Temps au bout duquel le signal se répète en un point. \\
\hline
Fréquence $f$ & Nombre de répétitions par seconde : $f=1/T$. \\
\hline
Longueur d'onde $\lambda$ & Distance entre deux points en phase ; distance parcourue pendant une période. \\
\hline
Relations & $c=\dfrac{d}{\Delta t}$, $\lambda=cT$, $c=\lambda f$. \\
\hline
\end{tabularx}
\end{center}

\section{La lumière : images et couleurs}

\subsection{La lentille mince convergente}
Une lentille convergente permet de faire converger des rayons lumineux et de former des images.

\begin{definitionbox}
Une \textbf{lentille mince convergente} est un système optique centré qui transforme un faisceau incident parallèle à l'axe optique en un faisceau convergent vers un point appelé \textbf{foyer image} $F'$.
\end{definitionbox}

\paragraph{Notations importantes.}
\begin{itemize}
    \item $O$ : centre optique ;
    \item $F$ : foyer objet ;
    \item $F'$ : foyer image ;
    \item $f'$ : distance focale image, égale à $OF'$.
\end{itemize}

\begin{center}
\begin{tikzpicture}[scale=1]
    \draw[->] (-5,0) -- (5.8,0) node[right] {axe optique};
    \draw[very thick] (0,-2) -- (0,2);
    \draw[blue,fill=blue] (0,2) -- (-0.15,1.7) -- (0.15,1.7) -- cycle;
    \draw[blue,fill=blue] (0,-2) -- (-0.15,-1.7) -- (0.15,-1.7) -- cycle;
    \node[below] at (0,0) {$O$};
    \fill (-2,0) circle (1.2pt) node[below] {$F$};
    \fill (2,0) circle (1.2pt) node[below] {$F'$};
    \draw[red,thick,->] (-4,1.2) -- (0,1.2);
    \draw[red,thick,->] (0,1.2) -- (2,0);
    \draw[red,thick,dashed] (2,0) -- (4,-1.2);
    \draw[orange,thick,->] (-4,-0.8) -- (0,-0.8);
    \draw[orange,thick,->] (0,-0.8) -- (4,-0.8);
    \node at (3.7,1.2) {rayon passant par $F'$};
\end{tikzpicture}
\end{center}

\subsection{Construction graphique d'une image}
Pour construire l'image d'un objet avec une lentille convergente, on utilise des rayons particuliers.

\begin{propriete}
Pour une lentille mince convergente :
\begin{itemize}
    \item un rayon passant par le centre optique $O$ n'est pas dévié ;
    \item un rayon incident parallèle à l'axe émerge en passant par $F'$ ;
    \item un rayon passant par $F$ émerge parallèle à l'axe.
\end{itemize}
\end{propriete}

\subsection{Image réelle et image virtuelle}

\begin{definitionbox}
\begin{itemize}
    \item Une \textbf{image réelle} est obtenue par intersection réelle des rayons émergents ; elle peut être recueillie sur un écran.
    \item Une \textbf{image virtuelle} est obtenue par intersection des prolongements des rayons émergents ; elle ne peut pas être projetée sur un écran.
\end{itemize}
\end{definitionbox}

\begin{propriete}
Avec une lentille convergente :
\begin{itemize}
    \item si l'objet est placé au-delà du foyer, l'image peut être réelle et renversée ;
    \item si l'objet est placé entre $O$ et $F$, l'image est virtuelle, droite et agrandie.
\end{itemize}
\end{propriete}

\subsection{Grandeurs algébriques et relation de conjugaison}
On travaille sur l'axe optique avec des distances orientées. On adopte la convention usuelle de lycée : les distances algébriques peuvent être négatives ou positives selon la position des points.

\begin{propriete}
La relation de conjugaison d'une lentille mince convergente s'écrit
\[
\frac{1}{\overline{OA'}}-\frac{1}{\overline{OA}}=\frac{1}{\overline{OF'}}=\frac{1}{f'}.
\]
\end{propriete}

\begin{propriete}
Le grandissement vaut
\[
\gamma = \frac{\overline{A'B'}}{\overline{AB}} = \frac{\overline{OA'}}{\overline{OA}}.
\]
\end{propriete}

\paragraph{Sens physique du grandissement.}
\begin{itemize}
    \item Si $|\gamma|>1$, l'image est plus grande que l'objet.
    \item Si $|\gamma|<1$, l'image est plus petite.
    \item Si $\gamma>0$, l'image est droite.
    \item Si $\gamma<0$, l'image est renversée.
\end{itemize}

\begin{demonstration}
La relation de grandissement vient du théorème de Thalès appliqué dans les triangles construits par les rayons issus du sommet de l'objet. Comme un rayon passant par $O$ n'est pas dévié, on obtient la proportion entre taille et distance à la lentille.
\end{demonstration}

\begin{exemplebox}
On place un objet à \SI{30}{cm} d'une lentille convergente de distance focale \SI{10}{cm}. En notation algébrique, on prend
\[
\overline{OA}=-\SI{30}{cm},\qquad \overline{OF'}=+\SI{10}{cm}.
\]
La relation de conjugaison donne
\[
\frac{1}{\overline{OA'}}-\frac{1}{-30}=\frac{1}{10}
\]
soit
\[
\frac{1}{\overline{OA'}}+\frac{1}{30}=\frac{1}{10}
\Rightarrow \frac{1}{\overline{OA'}}=\frac{1}{10}-\frac{1}{30}=\frac{2}{30}=\frac{1}{15}.
\]
Donc
\[
\overline{OA'}=+\SI{15}{cm}.
\]
L'image est réelle. Le grandissement vaut
\[
\gamma=\frac{15}{-30}=-0,5.
\]
L'image est renversée et deux fois plus petite que l'objet.
\end{exemplebox}

\subsection{Mise au point et distance focale}
Faire la mise au point consiste à obtenir une image nette sur le capteur ou sur l'écran. On peut y parvenir :
\begin{itemize}
    \item en déplaçant la lentille ou l'écran ;
    \item en changeant la distance focale dans certains systèmes optiques modernes.
\end{itemize}

\begin{ouverture}
En Terminale et dans les études supérieures, on utilisera plus souvent les modèles de formation d'image en combinant plusieurs systèmes optiques. La logique reste pourtant la même : un modèle géométrique simple pour prédire position et taille de l'image.
\end{ouverture}

\section{Couleurs et vision}

\subsection{Pourquoi voit-on des couleurs ?}
La couleur perçue dépend à la fois :
\begin{enumerate}
    \item de la lumière incidente ;
    \item de ce que l'objet absorbe, diffuse ou transmet ;
    \item de la manière dont notre œil et notre cerveau interprètent la lumière reçue.
\end{enumerate}

\begin{definitionbox}
\begin{itemize}
    \item \textbf{Absorption} : une partie de la lumière est captée par le matériau.
    \item \textbf{Transmission} : la lumière traverse le matériau.
    \item \textbf{Diffusion} : la lumière est renvoyée dans différentes directions.
\end{itemize}
\end{definitionbox}

\paragraph{Exemple fondamental.}
Un objet rouge éclairé en lumière blanche diffuse principalement le rouge et absorbe une grande partie des autres composantes.

\subsection{Synthèse additive}

\begin{definitionbox}
La \textbf{synthèse additive} consiste à superposer des lumières colorées. Les couleurs primaires additives sont \textbf{rouge, vert, bleu} (RVB).
\end{definitionbox}

\begin{center}
\begin{tikzpicture}[scale=1]
    \fill[red,opacity=0.45] (-0.9,0) circle (1.5);
    \fill[green!70!black,opacity=0.45] (0.9,0) circle (1.5);
    \fill[blue,opacity=0.45] (0,1.1) circle (1.5);
    \node at (-1.8,-1.2) {Rouge};
    \node at (1.8,-1.2) {Vert};
    \node at (0,2.45) {Bleu};
    \node at (0,-0.2) {Jaune};
    \node at (-0.45,0.95) {Magenta};
    \node at (0.45,0.95) {Cyan};
    \node[font=\bfseries] at (0,0.45) {Blanc};
\end{tikzpicture}
\end{center}

\begin{propriete}
En synthèse additive :
\begin{itemize}
    \item rouge + vert = jaune ;
    \item rouge + bleu = magenta ;
    \item vert + bleu = cyan ;
    \item rouge + vert + bleu = blanc.
\end{itemize}
\end{propriete}

\subsection{Synthèse soustractive}

\begin{definitionbox}
La \textbf{synthèse soustractive} correspond au mélange de pigments, d'encres ou à l'utilisation de filtres colorés qui retirent certaines composantes de la lumière. Les couleurs primaires soustractives sont \textbf{cyan, magenta, jaune} (CMJ).
\end{definitionbox}

\paragraph{Idée simple.}
Un filtre ne \og rajoute \fg{} pas une couleur : il enlève certaines radiations lumineuses. C'est pour cela que l'on parle de synthèse \emph{soustractive}.

\subsection{Couleurs complémentaires}

\begin{definitionbox}
Deux couleurs sont \textbf{complémentaires} lorsqu'en synthèse additive leur superposition donne du blanc.
\end{definitionbox}

\begin{propriete}
Les couples complémentaires usuels sont :
\begin{itemize}
    \item rouge / cyan ;
    \item vert / magenta ;
    \item bleu / jaune.
\end{itemize}
\end{propriete}

\subsection{Vision des couleurs et trichromie}
L'œil humain possède trois types de cônes sensibles à des zones différentes du spectre visible. La vision colorée repose donc sur une \textbf{trichromie}. Cela explique qu'on puisse reproduire beaucoup de couleurs à partir de trois lumières bien choisies.

\begin{ouverture}
Cette idée est fondamentale pour comprendre les écrans, capteurs photo, téléviseurs, vidéoprojecteurs et l'imagerie numérique. En Terminale, on reviendra sur des descriptions plus fines de la lumière et de son interaction avec la matière.
\end{ouverture}

\section{Modèles ondulatoire et particulaire de la lumière}

\subsection{La lumière comme onde électromagnétique}

\begin{definitionbox}
La lumière est une onde électromagnétique. Contrairement aux ondes mécaniques, elle peut se propager dans le vide.
\end{definitionbox}

\begin{propriete}
Dans le vide, la célérité de la lumière vaut
\[
c = 3{,}00\times 10^8\ \si{m.s^{-1}}.
\]
La relation entre célérité, longueur d'onde et fréquence est
\[
\lambda = \frac{c}{\nu}
\qquad \text{ou} \qquad c=\lambda\nu.
\]
\end{propriete}

\paragraph{Remarque de vocabulaire.}
Pour la lumière, on note souvent la fréquence $\nu$ au lieu de $f$.

\subsection{Domaines des ondes électromagnétiques}
Le visible n'est qu'une petite partie du spectre électromagnétique.

\begin{center}
\renewcommand{\arraystretch}{1.25}
\begin{tabularx}{\linewidth}{|>{\bfseries}m{4cm}|X|X|}
\hline
Domaine & Ordre de grandeur de $\lambda$ & Exemples d'usage \\
\hline
Ondes radio & très grandes longueurs d'onde & radio, communication \\
\hline
Micro-ondes & $\sim \SI{}{cm}$ à \SI{}{mm} & wifi, four micro-ondes, radar \\
\hline
Infrarouge & au-delà du rouge visible & thermique, télécommandes \\
\hline
Visible & environ de \SI{400}{nm} à \SI{800}{nm} & vision humaine \\
\hline
Ultraviolet & plus petit que le visible & bronzage, fluorescence \\
\hline
Rayons X & très petites longueurs d'onde & imagerie médicale \\
\hline
Rayons gamma & encore plus énergétiques & physique nucléaire, astrophysique \\
\hline
\end{tabularx}
\end{center}

\begin{attentionbox}
Plus la fréquence est grande, plus la longueur d'onde est petite. Il ne faut pas inverser cette idée.
\end{attentionbox}

\subsection{La lumière comme ensemble de photons}

\begin{definitionbox}
Le \textbf{photon} est le quantum d'énergie associé au rayonnement lumineux.
\end{definitionbox}

\begin{propriete}
L'énergie d'un photon vaut
\[
E = h\nu = \frac{hc}{\lambda}
\]
où $h=6{,}63\times 10^{-34}\ \si{J.s}$ est la constante de Planck.
\end{propriete}

\begin{demonstration}
Cette relation est un résultat expérimental fondamental de la physique quantique. Elle relie la description ondulatoire (fréquence, longueur d'onde) à la description particulaire (quanta d'énergie). Plus la fréquence est élevée, plus chaque photon transporte d'énergie.
\end{demonstration}

\begin{exemplebox}
Une lumière bleue a une fréquence plus grande qu'une lumière rouge. Donc ses photons sont plus énergétiques.
\end{exemplebox}

\subsection{Interaction lumière-matière et niveaux d'énergie}
Les atomes et molécules ne peuvent pas posséder n'importe quelle énergie : leurs niveaux d'énergie sont \textbf{quantifiés}.

\begin{definitionbox}
Dire que les niveaux d'énergie sont quantifiés signifie que l'atome ne peut exister que dans certaines valeurs d'énergie bien déterminées.
\end{definitionbox}

\begin{propriete}
Lorsqu'un atome passe d'un niveau d'énergie élevé à un niveau plus bas, il peut émettre un photon d'énergie
\[
\Delta E = h\nu.
\]
Lorsqu'il absorbe un photon, celui-ci doit avoir exactement l'énergie correspondant à l'écart entre deux niveaux.
\end{propriete}

\begin{center}
\begin{tikzpicture}[scale=1]
    \draw[thick] (0,0) -- (2.7,0) node[right] {$E_1$};
    \draw[thick] (0,1.2) -- (2.7,1.2) node[right] {$E_2$};
    \draw[thick] (0,2.5) -- (2.7,2.5) node[right] {$E_3$};
    \draw[->,red,very thick] (1.2,2.5) -- (1.2,1.2);
    \node[red,right] at (1.2,1.85) {émission};
    \draw[->,blue,very thick] (2,0) -- (2,1.2);
    \node[blue,right] at (2,0.6) {absorption};
\end{tikzpicture}
\end{center}

\paragraph{Conséquence importante.}
Les spectres d'émission et d'absorption ne sont pas continus pour les atomes isolés : on observe des \textbf{raies}. Chaque raie correspond à une transition entre niveaux d'énergie.

\begin{ouverture}
En Terminale, cette idée sera reliée plus loin à la quantification, aux spectres atomiques, à la structure de la matière et à des phénomènes comme l'effet photoélectrique.
\end{ouverture}

\section{Démonstrations et raisonnements à savoir refaire}

\subsection{Relation $\lambda=cT$}
\begin{demonstration}
Pendant la durée d'une période $T$, le motif élémentaire de l'onde s'est reproduit. La perturbation a donc avancé exactement d'une longueur d'onde $\lambda$. Par définition de la célérité,
\[
c=\frac{\lambda}{T} \quad \Longrightarrow \quad \lambda=cT.
\]
\end{demonstration}

\subsection{Relation $c=\lambda\nu$}
\begin{demonstration}
On a $\nu=1/T$. En remplaçant dans $\lambda=cT$, on obtient
\[
\lambda=c\times \frac{1}{\nu}
\quad \Longrightarrow \quad c=\lambda\nu.
\]
\end{demonstration}

\subsection{Sens du signe du grandissement}
\begin{demonstration}
Le grandissement compare algébriquement la taille orientée de l'image à celle de l'objet. Si le rapport est négatif, cela signifie que les segments sont orientés en sens opposés : l'image est renversée. S'il est positif, l'image est droite.
\end{demonstration}

\subsection{Lien entre fréquence et énergie du photon}
\begin{demonstration}
La relation fondamentale $E=h\nu$ montre que l'énergie d'un photon est proportionnelle à la fréquence. Ainsi, doubler la fréquence double l'énergie transportée par un photon. Comme $\lambda=c/\nu$, les plus petites longueurs d'onde correspondent donc aux photons les plus énergétiques.
\end{demonstration}

\section{Méthodes-types}

\subsection{Déterminer une célérité}
\begin{methode}
\begin{enumerate}
    \item Identifier la distance parcourue par la perturbation.
    \item Identifier la durée correspondante.
    \item Utiliser $c=d/\Delta t$.
    \item Vérifier les unités : mètre et seconde donnent des m\,s$^{-1}$.
\end{enumerate}
\end{methode}

\subsection{Déterminer une longueur d'onde}
\begin{methode}
Deux stratégies sont possibles.
\begin{itemize}
    \item \textbf{À partir d'un graphique spatial} : lire directement la distance entre deux crêtes successives.
    \item \textbf{À partir de $c$ et $f$} : utiliser $\lambda=c/f$.
\end{itemize}
\end{methode}

\subsection{Exploiter une lentille}
\begin{methode}
\begin{enumerate}
    \item Faire un schéma propre et placer l'objet, la lentille, $F$ et $F'$.
    \item Identifier si l'image attendue est réelle ou virtuelle selon la position de l'objet.
    \item Utiliser la relation de conjugaison.
    \item Utiliser le grandissement pour la taille et le sens de l'image.
\end{enumerate}
\end{methode}

\subsection{Interpréter une couleur}
\begin{methode}
Toujours se poser trois questions :
\begin{enumerate}
    \item Quelle est la lumière incidente ?
    \item Quelles composantes l'objet absorbe-t-il ?
    \item Quelles composantes diffuse-t-il ou transmet-il vers l'œil ?
\end{enumerate}
\end{methode}

\section{Exercices d'application}

\subsection*{Exercice 1 -- Propagation sur une corde}
Une impulsion se propage le long d'une corde. Elle parcourt \SI{1,8}{m} en \SI{0,60}{s}.
\begin{enumerate}[label=\alph*)]
    \item Calculer la célérité de l'onde.
    \item Que devient cette célérité si la corde est davantage tendue ?
\end{enumerate}

\textbf{Correction.}
\begin{enumerate}[label=\alph*)]
    \item
    \[
    c=\frac{1,8}{0,60}=\SI{3,0}{m.s^{-1}}.
    \]
    \item En général, si la corde est davantage tendue, l'onde se propage plus vite : la célérité augmente.
\end{enumerate}

\subsection*{Exercice 2 -- Longueur d'onde sonore}
Un diapason émet un son de fréquence \SI{512}{Hz}. La célérité du son dans l'air vaut \SI{340}{m.s^{-1}}.
\begin{enumerate}[label=\alph*)]
    \item Calculer la période du signal.
    \item Calculer la longueur d'onde.
\end{enumerate}

\textbf{Correction.}
\begin{enumerate}[label=\alph*)]
    \item
    \[
    T=\frac{1}{f}=\frac{1}{512}\approx 1{,}95\times 10^{-3}\ \si{s}.
    \]
    \item
    \[
    \lambda=\frac{c}{f}=\frac{340}{512}\approx \SI{0,664}{m}.
    \]
\end{enumerate}

\subsection*{Exercice 3 -- Lentille convergente}
On considère une lentille convergente de distance focale \SI{12}{cm}. Un objet de hauteur \SI{2,0}{cm} est placé à \SI{18}{cm} de la lentille.
\begin{enumerate}[label=\alph*)]
    \item Déterminer la position de l'image.
    \item Déterminer sa taille et préciser si elle est droite ou renversée.
\end{enumerate}

\textbf{Correction.}
On prend \(\overline{OA}=-18\ \si{cm}\) et \(\overline{OF'}=+12\ \si{cm}\).
\[
\frac{1}{\overline{OA'}}-\frac{1}{-18}=\frac{1}{12}
\Rightarrow \frac{1}{\overline{OA'}}+\frac{1}{18}=\frac{1}{12}
\]
\[
\frac{1}{\overline{OA'}}=\frac{1}{12}-\frac{1}{18}=\frac{1}{36}
\Rightarrow \overline{OA'}=+36\ \si{cm}.
\]
L'image est réelle. Le grandissement vaut
\[
\gamma=\frac{36}{-18}=-2.
\]
Donc
\[
A'B'=\gamma\times AB=-2\times 2{,}0=-4{,}0\ \si{cm}.
\]
L'image mesure \SI{4,0}{cm} et elle est renversée.

\subsection*{Exercice 4 -- Couleur d'un objet}
Une fleur rouge est éclairée successivement en lumière blanche, en lumière rouge puis en lumière verte. Préciser sa couleur apparente dans chaque cas.

\textbf{Correction.}
\begin{itemize}
    \item En lumière blanche : elle apparaît rouge, car elle diffuse principalement le rouge.
    \item En lumière rouge : elle apparaît rouge, car la lumière incidente contient du rouge.
    \item En lumière verte : elle apparaît noire ou très sombre, car elle absorbe l'essentiel de la lumière verte et ne peut pas diffuser du rouge absent de l'éclairage.
\end{itemize}

\subsection*{Exercice 5 -- Photon et longueur d'onde}
Entre une lumière rouge et une lumière violette, laquelle transporte les photons les plus énergétiques ? Justifier.

\textbf{Correction.}
L'énergie d'un photon vaut $E=h\nu=hc/\lambda$. La lumière violette a une longueur d'onde plus petite et une fréquence plus grande que la lumière rouge. Ses photons sont donc plus énergétiques.

\section{Vrai ou faux commenté}

\begin{enumerate}[label=\textbf{\arabic*.}]
    \item \textbf{Une onde mécanique peut se propager dans le vide.} \\
    \textbf{Faux.} Une onde mécanique nécessite un milieu matériel.

    \item \textbf{La célérité d'une onde et la vitesse d'un point du milieu sont toujours la même chose.} \\
    \textbf{Faux.} La célérité décrit la propagation de la perturbation ; les points du milieu ont en général un mouvement local différent.

    \item \textbf{La longueur d'onde est la distance parcourue pendant une période.} \\
    \textbf{Vrai.} C'est même une définition essentielle pour une onde périodique.

    \item \textbf{Si la fréquence augmente dans un même milieu, la longueur d'onde augmente aussi.} \\
    \textbf{Faux.} Comme $\lambda=c/f$ dans un milieu donné, si $f$ augmente alors $\lambda$ diminue.

    \item \textbf{Une image virtuelle peut être projetée sur un écran.} \\
    \textbf{Faux.} Seule une image réelle peut être recueillie sur un écran.

    \item \textbf{Un grandissement négatif signifie que l'image est renversée.} \\
    \textbf{Vrai.}

    \item \textbf{Un filtre rouge ajoute du rouge à la lumière blanche.} \\
    \textbf{Faux.} Il transmet surtout le rouge et absorbe une partie des autres couleurs ; il ne crée pas de lumière supplémentaire.

    \item \textbf{Le visible représente tout le spectre électromagnétique.} \\
    \textbf{Faux.} Le visible n'est qu'une petite bande du spectre.

    \item \textbf{Plus la fréquence de la lumière est grande, plus l'énergie du photon est grande.} \\
    \textbf{Vrai, car $E=h\nu$.}

    \item \textbf{Les spectres de raies montrent que les niveaux d'énergie sont quantifiés.} \\
    \textbf{Vrai.} Les transitions ne peuvent se faire qu'entre certaines valeurs d'énergie.
\end{enumerate}

\section{Questions de réflexion pour aller plus loin}
\begin{enumerate}[label=\arabic*.]
    \item Pourquoi entend-on parfois l'orage après avoir vu l'éclair ?
    \item Pourquoi un écran de téléphone peut-il fabriquer beaucoup de couleurs à partir de seulement trois sous-pixels ?
    \item Pourquoi une lumière bleue n'a-t-elle pas la même énergie par photon qu'une lumière rouge ?
    \item Pourquoi une image peut-elle être nette pour un objet et floue pour un autre ?
    \item Pourquoi le modèle ondulatoire et le modèle particulaire de la lumière sont-ils tous les deux utiles ?
\end{enumerate}

\section{Bilan final à connaître par cœur}
\begin{itemize}
    \item Une onde mécanique progressive est la propagation d'une perturbation dans un milieu matériel.
    \item La célérité vérifie $c=d/\Delta t$.
    \item Pour une onde périodique, $\lambda=cT$ et $c=\lambda f$.
    \item Une lentille convergente peut former une image réelle ou virtuelle selon la position de l'objet.
    \item Les relations essentielles en optique sont la relation de conjugaison et le grandissement.
    \item La couleur perçue dépend de la lumière incidente et de l'interaction avec l'objet.
    \item La synthèse additive concerne les lumières ; la synthèse soustractive concerne filtres et pigments.
    \item La lumière est une onde électromagnétique, mais aussi un flux de photons.
    \item L'énergie d'un photon vaut $E=h\nu=hc/\lambda$.
    \item Les niveaux d'énergie des atomes sont quantifiés.
\end{itemize}

\section*{Petit pont vers la Terminale}
Ce chapitre est une base. En Terminale, tu retrouveras plusieurs prolongements naturels :
\begin{itemize}
    \item description plus riche des ondes et signaux périodiques ;
    \item phénomènes d'interférences, de diffraction et parfois d'effet Doppler ;
    \item spectres et interaction lumière-matière plus approfondis ;
    \item outils mathématiques plus sûrs pour manipuler fonctions périodiques et modèles sinusoïdaux ;
    \item traitement expérimental plus poussé des signaux à l'aide d'interfaces et de logiciels.
\end{itemize}
La meilleure préparation consiste donc à être parfaitement à l'aise avec :
\[
 c=\frac{d}{\Delta t}, \qquad \lambda=cT, \qquad c=\lambda f, \qquad E=h\nu=\frac{hc}{\lambda}.
\]

\end{document}
